\documentclass[11pt]{article}
% Shared preamble for the LSM crash-course series.
\usepackage[a4paper,margin=0.9in]{geometry}
\usepackage{amsmath,amssymb,amsthm,mathtools}
\usepackage[dvipsnames]{xcolor}
\usepackage{enumitem}
\usepackage{booktabs}
\usepackage{longtable}
\usepackage{array}
\usepackage{tcolorbox}
\usepackage{fancyhdr}
\usepackage[colorlinks=true,linkcolor=Blue]{hyperref}
\tcbuselibrary{skins,breakable}

\setlist{itemsep=3pt,topsep=4pt}
\pagestyle{fancy}\fancyhf{}
\rhead{\small Linear Statistical Models, B.Stat.\ 3rd yr}
\cfoot{\small\thepage}

\newcommand{\E}{\mathbb{E}}
\newcommand{\Var}{\operatorname{Var}}
\newcommand{\Cov}{\operatorname{Cov}}
\newcommand{\Prob}{\mathbb{P}}
\newcommand{\R}{\mathbb{R}}
\newcommand{\tr}{\operatorname{tr}}
\newcommand{\rank}{\operatorname{rank}}
\newcommand{\Cs}{\mathcal{C}}
\newcommand{\N}{\mathcal{N}}
\newcommand{\Ybar}{\overline{Y}}
\newcommand{\Xbar}{\overline{X}}
\newcommand{\xbar}{\bar{x}}
\newcommand{\iid}{\overset{\text{iid}}{\sim}}
\newcommand{\indep}{\overset{\text{indep}}{\sim}}
\newcommand{\dto}{\overset{d}{\longrightarrow}}
\newcommand{\dequal}{\overset{d}{=}}
\newcommand{\bY}{\mathbf{Y}}
\newcommand{\bX}{\mathbf{X}}
\newcommand{\bZ}{\mathbf{Z}}
\newcommand{\bW}{\mathbf{W}}
\newcommand{\bb}{\boldsymbol{\beta}}
\newcommand{\bmu}{\boldsymbol{\mu}}
\newcommand{\bep}{\boldsymbol{\varepsilon}}
\newcommand{\bal}{\boldsymbol{\alpha}}
\newcommand{\one}{\mathbf{1}}
\newcommand{\zero}{\mathbf{0}}
\newcommand{\I}{\mathbb{I}}
\newcommand{\prf}{\smallskip\noindent\textbf{Proof.}\ }
\renewcommand{\qed}{\hfill$\blacksquare$}

\newtcolorbox{idea}[1]{colback=Blue!4,colframe=Blue!55,title={#1},
  fonttitle=\bfseries,breakable}
\newtcolorbox{confusion}[1]{colback=BrickRed!5,colframe=BrickRed!60,
  title={Common confusion: #1},fonttitle=\bfseries,breakable}
\newtcolorbox{recipe}[1]{colback=ForestGreen!5,colframe=ForestGreen!55,title={#1},
  fonttitle=\bfseries,breakable}
\newtcolorbox{exam}[1]{colback=Orange!6,colframe=Orange!70,title={#1},
  fonttitle=\bfseries,breakable}

\lhead{\small Previous-Year Papers --- Full Solutions}

\newcommand{\wh}[1]{\widehat{#1}}
\newcommand{\Pj}{\mathbf{P}}
\newcommand{\Nsp}{\mathcal{N}}
\newcommand{\SSE}{\mathrm{SSE}}
\newcommand{\MSE}{\mathrm{MSE}}
\newcommand{\Bt}{\mathrm{Beta}}
\newcommand{\ind}{\perp\!\!\!\perp}

\newcounter{qq}[section]
\newcommand{\qhead}[2]{\medskip\refstepcounter{qq}%
  \noindent\textbf{\color{Blue}Question \thesection.\theqq.}\hfill\textnormal{[\textbf{#1} marks]}%
  \ifx\relax#2\relax\else\par\smallskip\fi}
\newcommand{\sol}{\smallskip\noindent\textbf{\color{ForestGreen}Solution.}\ }
\newcommand{\pt}[1]{\smallskip\noindent\textbf{(#1)}\ }
\newtcolorbox{stmt}{colback=black!3,colframe=black!35,breakable,boxrule=0.4pt,
  left=4pt,right=4pt,top=3pt,bottom=3pt}

\begin{document}

\begin{center}
{\LARGE\bfseries Previous-Year Mid-Semester Papers}\\[4pt]
{\large Complete Worked Solutions}\\[3pt]
{\small Linear Statistical Models, B.Stat.\ 3rd year --- 2024--25 and 2025--26}
\end{center}
\vspace{2pt}\hrule\vspace{10pt}

\noindent Both papers are worth 40 marks. Every part is solved in full, with the
alternative route given wherever the professor's phrasing admits one. Section~3 adds five
worked practice questions in the same pure linear-algebra style, and Section~4 is a coverage
audit against the crash courses, recording where each result now lives.

\tableofcontents
\vspace{6pt}\hrule\vspace{8pt}

%==================================================================
\section{Mid-Semester Examination, 13 September 2024}

\qhead{10}{}
\begin{stmt}
$\mathbf{P}$ and $\mathbf{Q}$ are $n\times n$ symmetric orthogonal projection matrices.
Show that $\mathbf{P}+\mathbf{Q}$ is an orthogonal projection matrix if and only if
$\mathbf{PQ}=\mathbf{O}_n$. In that case show also that
$\rank(\mathbf{P})+\rank(\mathbf{Q})=\rank(\mathbf{P}+\mathbf{Q})$.
\end{stmt}

\sol Recall the characterisation: $M$ is an orthogonal projection matrix iff $M^T=M$ and
$M^2=M$. Since $\mathbf{P},\mathbf{Q}$ are symmetric, $\mathbf{P}+\mathbf{Q}$ is symmetric
automatically, so the whole question is about idempotence. Expand:
\[
  (\mathbf{P}+\mathbf{Q})^2=\mathbf{P}^2+\mathbf{PQ}+\mathbf{QP}+\mathbf{Q}^2
  =\mathbf{P}+\mathbf{Q}+(\mathbf{PQ}+\mathbf{QP}),
\]
so
\begin{equation}\label{eq:pq}
  (\mathbf{P}+\mathbf{Q})^2=\mathbf{P}+\mathbf{Q}
  \iff \mathbf{PQ}+\mathbf{QP}=\mathbf{O}.
\end{equation}

\pt{$\Leftarrow$} Suppose $\mathbf{PQ}=\mathbf{O}$. Since $\mathbf{P},\mathbf{Q}$ are
symmetric, $\mathbf{QP}=\mathbf{Q}^T\mathbf{P}^T=(\mathbf{PQ})^T=\mathbf{O}$ as well.
Hence $\mathbf{PQ}+\mathbf{QP}=\mathbf{O}$ and \eqref{eq:pq} gives idempotence.

\pt{$\Rightarrow$} Suppose $\mathbf{P}+\mathbf{Q}$ is an orthogonal projector, so by
\eqref{eq:pq}, $\mathbf{PQ}+\mathbf{QP}=\mathbf{O}$. Multiply this on the \emph{left} by
$\mathbf{P}$ and use $\mathbf{P}^2=\mathbf{P}$:
\[
  \mathbf{PQ}+\mathbf{PQP}=\mathbf{O}. \tag{i}
\]
Multiply it on the \emph{right} by $\mathbf{P}$:
\[
  \mathbf{PQP}+\mathbf{QP}=\mathbf{O}. \tag{ii}
\]
Subtracting (ii) from (i) gives $\mathbf{PQ}-\mathbf{QP}=\mathbf{O}$, i.e.\
$\mathbf{PQ}=\mathbf{QP}$. Substituting back into $\mathbf{PQ}+\mathbf{QP}=\mathbf{O}$
gives $2\mathbf{PQ}=\mathbf{O}$, hence $\mathbf{PQ}=\mathbf{O}$. \qed

\pt{Rank} Now assume $\mathbf{PQ}=\mathbf{O}$, so $\mathbf{P}+\mathbf{Q}$ is idempotent.
For an idempotent matrix all eigenvalues are $0$ or $1$, so the trace (the sum of the
eigenvalues) counts the eigenvalues equal to $1$, i.e.\ $\tr(M)=\rank(M)$. Applying this to
all three idempotent matrices and using linearity of the trace:
\[
  \rank(\mathbf{P}+\mathbf{Q})=\tr(\mathbf{P}+\mathbf{Q})=\tr(\mathbf{P})+\tr(\mathbf{Q})
  =\rank(\mathbf{P})+\rank(\mathbf{Q}). \qquad\square
\]

\begin{idea}{Alternative rank argument (no trace)}
Show $\Cs(\mathbf{P}+\mathbf{Q})=\Cs(\mathbf{P})\oplus\Cs(\mathbf{Q})$.
\begin{itemize}
  \item \emph{Orthogonality.} For $u=\mathbf{P}a$ and $v=\mathbf{Q}b$,
        $u^Tv=a^T\mathbf{PQ}b=0$. So $\Cs(\mathbf{P})\perp\Cs(\mathbf{Q})$, and in
        particular the sum is direct.
  \item ``$\subseteq$'': $(\mathbf{P}+\mathbf{Q})x=\mathbf{P}x+\mathbf{Q}x$.
  \item ``$\supseteq$'': $(\mathbf{P}+\mathbf{Q})\mathbf{P}a
        =\mathbf{P}^2a+\mathbf{QP}a=\mathbf{P}a$, so $\Cs(\mathbf{P})\subseteq\Cs(\mathbf{P}+\mathbf{Q})$;
        similarly for $\mathbf{Q}$.
\end{itemize}
Dimensions of a direct sum add, giving the result.
\end{idea}

%------------------------------------------------------------------
\qhead{12}{}
\begin{stmt}
Data $\bY=\{Y_{ij}: 1\le j\le m,\ i=1,2,3,4\}$, $m\ge2$, index $i$ the treatment and $j$ the
replicate, with
\[
  Y_{1j}=\beta_1+\varepsilon_{1j},\quad
  Y_{2j}=\beta_1+\beta_2+\varepsilon_{2j},\quad
  Y_{3j}=\beta_2+\beta_3+\varepsilon_{3j},\quad
  Y_{4j}=\beta_3+\beta_4+\varepsilon_{4j},
\]
$\varepsilon_{ij}\iid \N(0,\sigma^2)$.
\textbf{(a)} Write the model in matrix form.
\textbf{(b)} Are (i) $\beta_2$; (ii) $\beta_1-\beta_2$; (iii) $\beta_4$ estimable? For each
estimable one find the BLUE.
\textbf{(c)} Propose a test of $H_0:\beta_1=\beta_2$ against $H_A:\beta_1\ne\beta_2$.
\end{stmt}

\pt{a} Stack $\bY$ treatment-major,
$\bY=(Y_{11},\dots,Y_{1m},Y_{21},\dots,Y_{2m},Y_{31},\dots,Y_{4m})^T$, and put
$\bb=(\beta_1,\beta_2,\beta_3,\beta_4)^T$. Then $\bY=\bX\bb+\bep$ with
$\bep\sim\N_{4m}(\zero,\sigma^2I_{4m})$ and
\[
  \underset{4m\times4}{\bX}=
  \begin{pmatrix}
    \one_m & \zero & \zero & \zero\\
    \one_m & \one_m & \zero & \zero\\
    \zero & \one_m & \one_m & \zero\\
    \zero & \zero & \one_m & \one_m
  \end{pmatrix}
  = A\otimes\one_m,
  \qquad
  A=\begin{pmatrix}1&0&0&0\\1&1&0&0\\0&1&1&0\\0&0&1&1\end{pmatrix}.
\]

\pt{b} \textbf{First compute the rank --- do not assume it is deficient.} $A$ is lower
triangular with all diagonal entries equal to $1$, so $\det A=1$ and $A$ is invertible.
Hence $\rank(\bX)=\rank(A)=4=p$: the model has \textbf{full column rank}, so
$\bX^T\bX=m\,A^TA$ is non-singular, the model is identifiable, and

\begin{idea}{Every linear function of $\bb$ is estimable}
All three of $\beta_2$, $\beta_1-\beta_2$ and $\beta_4$ are estimable. This is the trap in
the question: the model \emph{looks} like an over-parametrised ANOVA, but it is not.
\end{idea}

Write $\mu_i:=\E(Y_{ij})$, so $\mu_1=\beta_1$, $\mu_2=\beta_1+\beta_2$,
$\mu_3=\beta_2+\beta_3$, $\mu_4=\beta_3+\beta_4$, i.e.\ $\bmu=A\bb$. Since
$\Cs(\bX)$ is exactly the span of the four treatment indicators (because $A$ is
invertible), the LSE of $\bmu$ is the vector of group means and, by forward substitution in
$A\bb=\bmu$,
\[
  \wh\bb=A^{-1}\begin{pmatrix}\Ybar_{1\cdot}\\ \Ybar_{2\cdot}\\ \Ybar_{3\cdot}\\ \Ybar_{4\cdot}\end{pmatrix}:
  \qquad
  \begin{aligned}
    \wh\beta_1&=\Ybar_{1\cdot}, &\qquad \wh\beta_3&=\Ybar_{3\cdot}-\Ybar_{2\cdot}+\Ybar_{1\cdot},\\
    \wh\beta_2&=\Ybar_{2\cdot}-\Ybar_{1\cdot}, &\qquad \wh\beta_4&=\Ybar_{4\cdot}-\Ybar_{3\cdot}+\Ybar_{2\cdot}-\Ybar_{1\cdot}.
  \end{aligned}
\]
(Check directly: $\hat\bb=(\bX^T\bX)^{-1}\bX^T\bY=(mA^TA)^{-1}(A^T\otimes\one_m^T)\bY
=(A^TA)^{-1}A^T\bar{\bY}=A^{-1}\bar{\bY}$, where $\bar\bY$ is the vector of group means.)

The $\Ybar_{i\cdot}$ are independent with variance $\sigma^2/m$, so for a coefficient vector
$c$ the variance is $\tfrac{\sigma^2}{m}\sum_ic_i^2$. By Gauss--Markov each of these is the
unique BLUE:

\begin{center}
\begin{tabular}{@{}llll@{}}
\toprule
& Parameter & BLUE & Variance\\
\midrule
(i) & $\beta_2$ & $\Ybar_{2\cdot}-\Ybar_{1\cdot}$ & $2\sigma^2/m$\\
(ii) & $\beta_1-\beta_2=2\mu_1-\mu_2$ & $2\Ybar_{1\cdot}-\Ybar_{2\cdot}$ & $5\sigma^2/m$\\
(iii) & $\beta_4$ & $\Ybar_{4\cdot}-\Ybar_{3\cdot}+\Ybar_{2\cdot}-\Ybar_{1\cdot}$ & $4\sigma^2/m$\\
\bottomrule
\end{tabular}
\end{center}

\pt{c} Write the hypothesis as $L^T\bb=\theta_0$ with $L^T=(1,-1,0,0)$, $\theta_0=0$,
so $r=\rank(L)=1$. Then
\[
  \wh\theta=\wh\beta_1-\wh\beta_2=2\Ybar_{1\cdot}-\Ybar_{2\cdot},
  \qquad
  \Var(\wh\theta)=\frac{5\sigma^2}{m}.
\]
Because $\Cs(\bX)$ is the span of the group indicators, the fitted values are
$\wh Y_{ij}=\Ybar_{i\cdot}$ and
\[
  \SSE=\sum_{i=1}^4\sum_{j=1}^m(Y_{ij}-\Ybar_{i\cdot})^2,
  \qquad
  \wh\sigma^2=\MSE=\frac{\SSE}{4m-4},
  \qquad \SSE\sim\sigma^2\chi^2_{4m-4}\ \text{independently of }\wh\bb.
\]
Hence the test statistic
\[
  \boxed{\ T=\frac{2\Ybar_{1\cdot}-\Ybar_{2\cdot}}{\wh\sigma\sqrt{5/m}}
  \ \overset{H_0}{\sim}\ t_{4m-4}\ },
  \qquad\text{reject }H_0\text{ if } |T|>t^{1-\alpha/2}_{4m-4}.
\]
Equivalently (and this is the ``alternative form of the likelihood ratio test'' the note
permits) use
\[
  F=T^2=\frac{W}{r}=\frac{(\SSE_{red}-\SSE_{full})/1}{\MSE}\ \overset{H_0}{\sim}\ F_{1,\,4m-4},
\]
where $\SSE_{red}$ is the residual sum of squares after fitting under the constraint
$\beta_1=\beta_2$.

%------------------------------------------------------------------
\qhead{6}{}
\begin{stmt}
$A$ is $m\times m$ of the form $A=\begin{bsmallmatrix}A_{11}&A_{12}\\A_{21}&A_{22}\end{bsmallmatrix}$
with $A_{11}$ of size $r\times r$, $r<m$, and $\rank(A)=\rank(A_{11})=r$. Show that
$G=\begin{bsmallmatrix}A_{11}^{-1}&0\\0&0\end{bsmallmatrix}$ is a generalised inverse of $A$.
\end{stmt}

\sol $A_{11}$ is $r\times r$ with $\rank(A_{11})=r$, so $A_{11}^{-1}$ exists and $G$ is well
defined. We must verify $AGA=A$. First,
\[
  AG=\begin{pmatrix}A_{11}&A_{12}\\A_{21}&A_{22}\end{pmatrix}
     \begin{pmatrix}A_{11}^{-1}&0\\0&0\end{pmatrix}
   =\begin{pmatrix}I_r&0\\A_{21}A_{11}^{-1}&0\end{pmatrix},
\]
and therefore
\begin{equation}\label{eq:aga}
  AGA=\begin{pmatrix}I_r&0\\A_{21}A_{11}^{-1}&0\end{pmatrix}
      \begin{pmatrix}A_{11}&A_{12}\\A_{21}&A_{22}\end{pmatrix}
    =\begin{pmatrix}A_{11}&A_{12}\\A_{21}&A_{21}A_{11}^{-1}A_{12}\end{pmatrix}.
\end{equation}
Comparing with $A$, it remains \emph{only} to show
\begin{equation}\label{eq:schur}
  A_{22}=A_{21}A_{11}^{-1}A_{12},
\end{equation}
i.e.\ that the Schur complement $A_{22}-A_{21}A_{11}^{-1}A_{12}$ vanishes.

\textbf{This is exactly what the rank hypothesis gives.} Write $A$ in terms of its two row
blocks, $A=\begin{bsmallmatrix}[A_{11}\ A_{12}]\\ [A_{21}\ A_{22}]\end{bsmallmatrix}$. The
top block $[A_{11}\ A_{12}]$ has $r$ rows and already has rank $r$ (since $A_{11}$ does), and
$\rank(A)=r$ by hypothesis. So the row space of $A$ is spanned by the first $r$ rows, and
every row of the bottom block is a linear combination of them: there is an
$(m-r)\times r$ matrix $M$ with
\[
  [A_{21}\ \ A_{22}]=M[A_{11}\ \ A_{12}]
  \quad\Longrightarrow\quad
  A_{21}=MA_{11}\ \ \text{and}\ \ A_{22}=MA_{12}.
\]
Since $A_{11}$ is invertible, the first equation gives $M=A_{21}A_{11}^{-1}$, and
substituting into the second gives \eqref{eq:schur}. With \eqref{eq:schur},
\eqref{eq:aga} reads $AGA=A$. \qed

\begin{confusion}{where the hypothesis $\rank(A)=r$ is used}
Only in \eqref{eq:schur}. If you drop it, $G$ still satisfies $AG$ idempotent, but $AGA$
differs from $A$ in the bottom-right block --- so $G$ would \emph{not} be a g-inverse. Say
explicitly where you use the rank condition; that is where the marks are.
\end{confusion}

%------------------------------------------------------------------
\qhead{12}{}
\begin{stmt}
$\bY=\bX\bb+\bep$, $\bX$ is $n\times p$ with $\rank(\bX)=p$, $\bep\sim\N_n(\zero,\sigma^2I_n)$.
$\wh\bb=(\wh\beta_1,\dots,\wh\beta_p)^T$ is the LSE and $\wh\sigma^2=\MSE$. Let
$\bX_{\cdot j}$ be the $j$-th column and $\bX_{-j}$ the $n\times(p-1)$ submatrix with the
$j$-th column removed.
\textbf{(a)} Show $\Var(\wh\beta_j)=\sigma^2/\bX_{\cdot j}^T(I-\Pj_{-j})\bX_{\cdot j}$.
\textbf{(b)} Show the $t$-statistic for $H_0:\beta_j=0$ can be written
$T_j=\frac{1}{\wh\sigma/\sigma}\big(\frac{\beta_j}{\sigma}\sqrt{\bX_{\cdot j}^T(I-\Pj_{-j})\bX_{\cdot j}}+W_j\big)$
with $W_j\sim\N(0,1)$ independent of $\wh\sigma^2$.
\textbf{(c)} If $\E(\bY)=\beta_{j_0}\bX_{\cdot j_0}$ and $\bX_{\cdot j_0}$ is nearly collinear
with the other columns, argue that the $t$-test in the full model is much less significant
than in the true model.
\end{stmt}

Throughout write $d_j:=\bX_{\cdot j}^T(I-\Pj_{-j})\bX_{\cdot j}=\|(I-\Pj_{-j})\bX_{\cdot j}\|^2$.

\pt{a} Apply Frisch--Waugh with $\bZ=\bX_{-j}$ (the other $p-1$ columns) and
$\bW=\bX_{\cdot j}$ (one column, coefficient $\beta_j$):
\[
  \wh\beta_j=\big(\bX_{\cdot j}^T(I-\Pj_{-j})\bX_{\cdot j}\big)^{-1}\bX_{\cdot j}^T(I-\Pj_{-j})\bY
  =\frac{1}{d_j}\,\bX_{\cdot j}^T(I-\Pj_{-j})\bY .
\]
So $\wh\beta_j=a^T\bY$ with $a=(I-\Pj_{-j})\bX_{\cdot j}/d_j$, and since
$\Var(\bY)=\sigma^2I$ and $I-\Pj_{-j}$ is symmetric idempotent,
\[
  \Var(\wh\beta_j)=\sigma^2a^Ta
  =\frac{\sigma^2}{d_j^2}\,\bX_{\cdot j}^T(I-\Pj_{-j})(I-\Pj_{-j})\bX_{\cdot j}
  =\frac{\sigma^2}{d_j^2}\cdot d_j=\frac{\sigma^2}{d_j}. \qquad\square
\]

\begin{idea}{Second route: Schur complement}
$\Var(\wh\bb)=\sigma^2(\bX^T\bX)^{-1}$, so $\Var(\wh\beta_j)=\sigma^2[(\bX^T\bX)^{-1}]_{jj}$.
Permute the $j$-th column to the end and use the partitioned inverse: the bottom-right entry
of $\begin{bsmallmatrix}\bX_{-j}^T\bX_{-j}&\bX_{-j}^T\bX_{\cdot j}\\
\bX_{\cdot j}^T\bX_{-j}&\bX_{\cdot j}^T\bX_{\cdot j}\end{bsmallmatrix}^{-1}$ is the
reciprocal of the Schur complement
$\bX_{\cdot j}^T\bX_{\cdot j}-\bX_{\cdot j}^T\bX_{-j}(\bX_{-j}^T\bX_{-j})^{-1}\bX_{-j}^T\bX_{\cdot j}
=\bX_{\cdot j}^T(I-\Pj_{-j})\bX_{\cdot j}=d_j$.
\end{idea}

\pt{b} From (a), $\wh\beta_j\sim\N(\beta_j,\sigma^2/d_j)$. Standardising,
\[
  W_j:=\frac{(\wh\beta_j-\beta_j)\sqrt{d_j}}{\sigma}\sim\N(0,1)
  \qquad\Longrightarrow\qquad
  \frac{\wh\beta_j\sqrt{d_j}}{\sigma}=\frac{\beta_j\sqrt{d_j}}{\sigma}+W_j .
\]
The $t$-statistic is $T_j=\wh\beta_j/s(\wh\beta_j)$ with $s(\wh\beta_j)=\wh\sigma/\sqrt{d_j}$,
so, inserting $\sigma/\sigma$,
\[
  T_j=\frac{\wh\beta_j\sqrt{d_j}}{\wh\sigma}
     =\frac{\sigma}{\wh\sigma}\cdot\frac{\wh\beta_j\sqrt{d_j}}{\sigma}
     =\frac{1}{\wh\sigma/\sigma}\left(\frac{\beta_j}{\sigma}\sqrt{d_j}+W_j\right).
\]
Finally $W_j$ is a function of $\wh\bb$ alone, and $\wh\bb\ind\SSE$ (shown by
$\E[(I-\Pj_\bX)\bep\,\bep^T\bX(\bX^T\bX)^{-1}]=\sigma^2(I-\Pj_\bX)\bX(\bX^T\bX)^{-1}=0$
plus joint normality), so $W_j\ind\wh\sigma^2$. \qed

Note that under $H_0:\beta_j=0$ this is $T_j=W_j/\sqrt{\wh\sigma^2/\sigma^2}$ with
$(n-p)\wh\sigma^2/\sigma^2\sim\chi^2_{n-p}$, i.e.\ exactly $t_{n-p}$.

\pt{c} Part (b) shows the whole signal in $T_{j_0}$ sits in the \textbf{non-centrality}
\[
  \frac{\beta_{j_0}}{\sigma}\sqrt{d_{j_0}},
  \qquad d_{j_0}=\|(I-\Pj_{-j_0})\bX_{\cdot j_0}\|^2 .
\]
\begin{itemize}
  \item \textbf{Full model.} ``$\bX_{\cdot j_0}$ nearly collinear with the other columns''
        means $\bX_{\cdot j_0}$ is close to $\Cs(\bX_{-j_0})$, so
        $(I-\Pj_{-j_0})\bX_{\cdot j_0}\approx\zero$ and $d_{j_0}\approx0$. Then the
        non-centrality is $\approx0$ and
        $T_{j_0}\approx W_{j_0}/(\wh\sigma/\sigma)$, which is (approximately) a
        \emph{central} $t_{n-p}$ --- the same distribution as under $H_0$. The test
        therefore has power close to its level: it will typically fail to reject even though
        $\beta_{j_0}\ne0$.
  \item \textbf{True model} $\bY\sim\N_n(\beta_{j_0}\bX_{\cdot j_0},\sigma^2I)$. Here there
        are no other columns to project out, so $d=\|\bX_{\cdot j_0}\|^2$, which is large.
        The non-centrality $\beta_{j_0}\|\bX_{\cdot j_0}\|/\sigma$ is large and the test is
        highly significant. (The error variance also gets $n-1$ degrees of freedom instead
        of $n-p$.)
\end{itemize}
Equivalently, $\Var(\wh\beta_{j_0})=\sigma^2/d_{j_0}$ is inflated by collinearity. Note
carefully that $\wh\beta_{j_0}$ remains \textbf{unbiased} in the full model --- collinearity
destroys precision and hence power, it does not introduce bias.

%==================================================================
\section{Mid-Semester Examination, 9 September 2025}

\qhead{5}{}
\begin{stmt}
$\bY\sim\N_n(\bX\bb,\sigma^2I_n)$, $\bb\in\R^p$, $\rank(\bX)=p$, $\bX=[\bZ:\bW]$ with $\bZ$
an $n\times q$ matrix, $1\le q<p$. Define
$S_1=\min_{\bb\in\R^p}\|\bY-\bX\bb\|_2^2$ and $S_0=\min_{\gamma\in\R^q}\|\bY-\bZ\gamma\|_2^2$.
Show $S_1/S_0\sim\Bt(k_1,k_2)$ with $k_1=(n-p)/2$ and $k_2=(p-q)/2$, provided
$\E(\bY)\in\Cs(\bZ)$.
\end{stmt}

\sol Both minima are squared lengths of residual vectors, i.e.\ quadratic forms in
orthogonal projections:
\[
  S_1=\|(I-\Pj_\bX)\bY\|^2=\bY^T(I-\Pj_\bX)\bY,
  \qquad
  S_0=\|(I-\Pj_\bZ)\bY\|^2=\bY^T(I-\Pj_\bZ)\bY.
\]
Write $S_0=S_1+(S_0-S_1)$ where, since $\Cs(\bZ)\subseteq\Cs(\bX)$,
\[
  S_0-S_1=\bY^T\big(\Pj_\bX-\Pj_\bZ\big)\bY=\bY^T\Pj_{\bW\mid\bZ}\bY .
\]
Put $A_1=I-\Pj_\bX$ and $A_2=\Pj_\bX-\Pj_\bZ$. Both are symmetric and idempotent
($A_2$ by the decomposition lemma $\Pj_\bX=\Pj_\bZ+\Pj_{\bW\mid\bZ}$), with
\[
  \rank(A_1)=n-\rank(\bX)=n-p,
  \qquad
  \rank(A_2)=\rank(\bX)-\rank(\bZ)=p-q.
\]

\textbf{Independence.} Using $\Pj_\bX\Pj_\bZ=\Pj_\bZ\Pj_\bX=\Pj_\bZ$,
\[
  A_1A_2=(I-\Pj_\bX)(\Pj_\bX-\Pj_\bZ)=\Pj_\bX-\Pj_\bZ-\Pj_\bX+\Pj_\bZ=O .
\]
Two quadratic forms in a normal vector with $A_1\Sigma A_2=\sigma^2A_1A_2=O$ are
independent.

\textbf{Both are central.} By hypothesis $\bmu:=\E(\bY)=\bX\bb\in\Cs(\bZ)$, so
$\Pj_\bZ\bmu=\bmu$ and also $\Pj_\bX\bmu=\bmu$. Hence
\[
  A_1\bmu=\bmu-\bmu=\zero, \qquad A_2\bmu=\bmu-\bmu=\zero,
\]
so both non-centrality parameters vanish and, by the quadratic-form theorem,
\[
  \frac{S_1}{\sigma^2}\sim\chi^2_{n-p},
  \qquad
  \frac{S_0-S_1}{\sigma^2}\sim\chi^2_{p-q},
  \qquad\text{independently.}
\]

\textbf{Beta step.} Write $A=S_1/\sigma^2$ and $B=(S_0-S_1)/\sigma^2$. Then
\[
  \frac{S_1}{S_0}=\frac{S_1}{S_1+(S_0-S_1)}=\frac{A}{A+B}.
\]
If $A\sim\chi^2_a=\text{Gamma}(a/2,2)$ and $B\sim\chi^2_b=\text{Gamma}(b/2,2)$ are
independent, the change of variables $(U,V)=\big(\tfrac{A}{A+B},A+B\big)$ has Jacobian $v$
and factorises the joint density into $u^{a/2-1}(1-u)^{b/2-1}$ times a
$\text{Gamma}(\tfrac{a+b}{2},2)$ density; hence $U\sim\Bt(a/2,b/2)$ independently of $V$.
With $a=n-p$ and $b=p-q$,
\[
  \boxed{\ \frac{S_1}{S_0}\sim\Bt\Big(\frac{n-p}{2},\ \frac{p-q}{2}\Big).\ }\qquad\square
\]

\begin{confusion}{which way up?}
$S_1\le S_0$ always (the bigger model fits at least as well), so the ratio $S_1/S_0$ lies in
$[0,1]$ --- as a Beta variate must. If you accidentally write $S_0/S_1$ you get a quantity
$\ge1$, which cannot be Beta. This is the fastest sanity check on the orientation, and it is
also why the first parameter is $(n-p)/2$ (the \emph{numerator's} degrees of freedom).
\end{confusion}

%------------------------------------------------------------------
\qhead{15}{}
\begin{stmt}
$Y_{ijk}=\alpha_i+\beta_j+\varepsilon_{ijk}$, $1\le i,j\le2$, $k=1,\dots,m$,
$\varepsilon_{ijk}\iid\N(0,\sigma^2)$, $m\ge2$.
\textbf{(a)} [3] Show the design matrix $\bX$ has rank 3.
\textbf{(b)} [4] Find the BLUE of $\theta=\bX^T\bmu$, where $\bmu=\E(\bY)$ has entries
$\mu_{ij}=\E(Y_{ijk})$.
\textbf{(c)} [8] For each of (i) $\alpha_1$; (ii) $\alpha_1-\alpha_2$;
(iii) $\alpha_1+\alpha_2+\beta_1+\beta_2$; (iv) $\alpha_1-\alpha_2+\beta_1-\beta_2$, check
estimability and, if estimable, find the BLUE. Justify.
\end{stmt}

Order the parameters as $\bb=(\alpha_1,\alpha_2,\beta_1,\beta_2)^T$ and the observations by
cell $(1,1),(1,2),(2,1),(2,2)$, each cell contributing $m$ identical rows. The four distinct
rows of $\bX$ are
\[
  (1,0,1,0),\quad (1,0,0,1),\quad (0,1,1,0),\quad (0,1,0,1),
\]
so $\bX=A\otimes\one_m$ with $A$ the $4\times4$ matrix of those rows.

\pt{a} \textbf{Upper bound.} Writing $c_1,\dots,c_4$ for the columns of $\bX$, both
$c_1+c_2$ and $c_3+c_4$ equal $\one_{4m}$ (every observation has exactly one $\alpha$ and
exactly one $\beta$). Hence
\[
  c_1+c_2-c_3-c_4=\zero \quad\Longrightarrow\quad \rank(\bX)\le3 .
\]
\textbf{Lower bound.} The rows $(1,0,1,0)$, $(1,0,0,1)$, $(0,1,1,0)$ are linearly
independent: if $a(1,0,1,0)+b(1,0,0,1)+c(0,1,1,0)=\zero$ then the second coordinate gives
$c=0$, the fourth gives $b=0$, and the first then gives $a=0$. Hence $\rank(\bX)\ge3$.
Therefore $\rank(\bX)=3$. \qed

Consequently $\dim\Nsp(\bX)=4-3=1$ and, from the dependency above,
\begin{equation}\label{eq:null}
  \Nsp(\bX)=\text{span}\{z\},\qquad z=(1,1,-1,-1)^T .
\end{equation}
(The model is not identifiable: replacing $\alpha_i\mapsto\alpha_i+c$,
$\beta_j\mapsto\beta_j-c$ leaves every $\mu_{ij}$ unchanged.)

\pt{b} Since $\bmu=\bX\bb$, we have $\theta=\bX^T\bmu=\bX^T\bX\bb$. Each coordinate of
$\theta$ is $a^T\bb$ with $a=\bX^T\bX e_k\in\Cs(\bX^T\bX)$, so
\textbf{$\theta$ is estimable without any condition}. Its BLUE is $a^T\wh\bb$ evaluated at
any solution of the normal equations, i.e.\ $\bX^T\bX\wh\bb=\bX^T\bY$, giving
\[
  \boxed{\ \wh\theta=\bX^T\bY\ }
  \qquad\text{(equivalently } \bX^T\wh\bmu=\bX^T\Pj_\bX\bY=(\Pj_\bX\bX)^T\bY=\bX^T\bY).
\]
Explicitly, $\bX^T\bY$ collects the row and column totals:
\[
  \wh\theta=\big(T_{1\cdot\cdot},\ T_{2\cdot\cdot},\ T_{\cdot1\cdot},\ T_{\cdot2\cdot}\big)^T,
  \qquad
  T_{i\cdot\cdot}=\sum_{j,k}Y_{ijk},\quad T_{\cdot j\cdot}=\sum_{i,k}Y_{ijk},
\]
with $\Var(\wh\theta)=\sigma^2\bX^T\bX
=\sigma^2m\begin{bsmallmatrix}2&0&1&1\\0&2&1&1\\1&1&2&0\\1&1&0&2\end{bsmallmatrix}$.

\pt{c} By \eqref{eq:null} and the criterion $a\in\Cs(\bX^T)=\Nsp(\bX)^\perp$,
\begin{equation}\label{eq:crit}
  a^T\bb \text{ is estimable} \iff a^Tz=0 \iff a_1+a_2=a_3+a_4 .
\end{equation}
For the BLUEs, note that in this balanced additive design the LS fitted cell means are
\[
  \wh\mu_{ij}=\Ybar_{i\cdot\cdot}+\Ybar_{\cdot j\cdot}-\Ybar_{\cdots},
\]
and any estimable function is a linear combination of the $\mu_{ij}$, whose BLUE is the same
combination of the $\wh\mu_{ij}$. Write $M_{ij}:=\Ybar_{ij\cdot}$ for the four cell means;
they are independent with variance $\sigma^2/m$.

\pt{i} $\alpha_1$: $a=(1,0,0,0)$, $a^Tz=1\ne0$. \textbf{Not estimable.} \\
Justification in words: for any $c$, the parameter vectors $\bb$ and
$\bb+cz$ give identical distributions for $\bY$, but different values of $\alpha_1$; so no
statistic can have expectation $\alpha_1$ for all $\bb$.

\pt{ii} $\alpha_1-\alpha_2$: $a=(1,-1,0,0)$, $a^Tz=1-1=0$. \textbf{Estimable.}
It equals $\mu_{1j}-\mu_{2j}$ for either $j$, so
\[
  \widehat{\alpha_1-\alpha_2}=\wh\mu_{11}-\wh\mu_{21}
  =\big(\Ybar_{1\cdot\cdot}+\Ybar_{\cdot1\cdot}-\Ybar_{\cdots}\big)-\big(\Ybar_{2\cdot\cdot}+\Ybar_{\cdot1\cdot}-\Ybar_{\cdots}\big)
  =\boxed{\Ybar_{1\cdot\cdot}-\Ybar_{2\cdot\cdot}} .
\]
$\Ybar_{1\cdot\cdot}$ and $\Ybar_{2\cdot\cdot}$ are means of $2m$ observations each, over
disjoint sets, hence independent with variance $\sigma^2/(2m)$; so
$\Var=\sigma^2/m$.

\pt{iii} $\alpha_1+\alpha_2+\beta_1+\beta_2$: $a=(1,1,1,1)$, $a^Tz=1+1-1-1=0$.
\textbf{Estimable.} It equals $(\alpha_1+\beta_1)+(\alpha_2+\beta_2)=\mu_{11}+\mu_{22}$
(and also $\mu_{12}+\mu_{21}$). Then
\[
  \wh\mu_{11}+\wh\mu_{22}
  =\big(\Ybar_{1\cdot\cdot}+\Ybar_{2\cdot\cdot}\big)+\big(\Ybar_{\cdot1\cdot}+\Ybar_{\cdot2\cdot}\big)-2\Ybar_{\cdots}
  =2\Ybar_{\cdots}+2\Ybar_{\cdots}-2\Ybar_{\cdots}
  =\boxed{2\Ybar_{\cdots}},
\]
using $\Ybar_{1\cdot\cdot}+\Ybar_{2\cdot\cdot}=\Ybar_{\cdot1\cdot}+\Ybar_{\cdot2\cdot}=2\Ybar_{\cdots}$
in the balanced design. $\Var=4\cdot\sigma^2/(4m)=\sigma^2/m$.

\pt{iv} $\alpha_1-\alpha_2+\beta_1-\beta_2$: $a=(1,-1,1,-1)$, $a^Tz=1-1-1+1=0$.
\textbf{Estimable.} It equals $(\alpha_1+\beta_1)-(\alpha_2+\beta_2)=\mu_{11}-\mu_{22}$, so
\[
  \wh\mu_{11}-\wh\mu_{22}
  =\big(\Ybar_{1\cdot\cdot}-\Ybar_{2\cdot\cdot}\big)+\big(\Ybar_{\cdot1\cdot}-\Ybar_{\cdot2\cdot}\big)
  =\boxed{\text{(row contrast)}+\text{(column contrast)}} .
\]
In terms of the cell means, the row contrast is $R=\tfrac12(M_{11}+M_{12}-M_{21}-M_{22})$ and
the column contrast is $C=\tfrac12(M_{11}-M_{12}+M_{21}-M_{22})$; each has variance
$\tfrac14\cdot4\cdot\tfrac{\sigma^2}{m}=\tfrac{\sigma^2}{m}$, and
\[
  \Cov(R,C)=\tfrac14\cdot\tfrac{\sigma^2}{m}\big[(1)(1)+(1)(-1)+(-1)(1)+(-1)(-1)\big]=0 ,
\]
so row and column contrasts are uncorrelated (orthogonality of the balanced design) and
$\Var=2\sigma^2/m$.

%------------------------------------------------------------------
\qhead{12}{}
\begin{stmt}
$Y_{ij}=\mu_i+\beta_1z_{ij}+\beta_2w_{ij}+\varepsilon_{ij}$, $j=1,\dots,n_i$, $i=1,\dots,r$,
$\varepsilon_{ij}\iid\N(0,\sigma^2)$, $n=\sum_in_i$. With
$\bar z_{i\cdot}=\tfrac1{n_i}\sum_jz_{ij}$ etc.,
\[
  S_{zz}=\sum_i\sum_j(z_{ij}-\bar z_{i\cdot})^2,\quad
  S_{zw}=\sum_i\sum_j(z_{ij}-\bar z_{i\cdot})(w_{ij}-\bar w_{i\cdot}),\quad
  S_{ww}=\sum_i\sum_j(w_{ij}-\bar w_{i\cdot})^2,
\]
\[
  S_{zy}=\sum_i\sum_j(z_{ij}-\bar z_{i\cdot})(Y_{ij}-\Ybar_{i\cdot}),\qquad
  S_{wy}=\sum_i\sum_j(w_{ij}-\bar w_{i\cdot})(Y_{ij}-\Ybar_{i\cdot}),
\]
and $\wh\sigma^2=\MSE$.
\textbf{(a)} [4] Show $\wh\bb=\begin{bsmallmatrix}S_{zz}&S_{zw}\\S_{zw}&S_{ww}\end{bsmallmatrix}^{-1}\begin{bsmallmatrix}S_{zy}\\S_{wy}\end{bsmallmatrix}$ for $\bb=(\beta_1,\beta_2)^T$.
\textbf{(b)} [4] Show the LRT of $H_0:\beta_2=0$ rejects for large
$U=\big(S_{ww}-S_{zw}^2/S_{zz}\big)\wh\beta_2^{\,2}/\wh\sigma^2$.
\textbf{(c)} [4] Show the LRT of $H_0:\beta_1=\beta_2=0$ rejects for large
$V=\tfrac1{\wh\sigma^2}\wh\bb^T\begin{bsmallmatrix}S_{zz}&S_{zw}\\S_{zw}&S_{ww}\end{bsmallmatrix}\wh\bb$.
\end{stmt}

Write the model as $\bY=\bZ\gamma+\bW\bb+\bep$, where $\bZ$ is the $n\times r$ matrix of
treatment indicators, $\gamma=(\mu_1,\dots,\mu_r)^T$, and $\bW=[\,z:w\,]$ is $n\times2$. The
full model has $p=r+2$ parameters, so $\MSE$ has $n-r-2$ degrees of freedom.

\begin{idea}{The one observation that makes all three parts easy}
$\Pj_\bZ$ replaces each observation by its \emph{treatment mean}, so $I-\Pj_\bZ$
\textbf{centres within treatment}:
\[
  \big[(I-\Pj_\bZ)v\big]_{ij}=v_{ij}-\bar v_{i\cdot}.
\]
Every $S$ in the question is therefore an inner product of two such centred vectors, and
since $I-\Pj_\bZ$ is symmetric idempotent,
\[
  u^T(I-\Pj_\bZ)v=\big((I-\Pj_\bZ)u\big)^T\big((I-\Pj_\bZ)v\big)
  =\sum_i\sum_j(u_{ij}-\bar u_{i\cdot})(v_{ij}-\bar v_{i\cdot}).
\]
\end{idea}

\pt{a} By Frisch--Waugh applied with $\bZ$ partialled out,
\[
  \wh\bb=\big(\bW^T(I-\Pj_\bZ)\bW\big)^{-1}\bW^T(I-\Pj_\bZ)\bY .
\]
By the boxed identity,
\[
  \bW^T(I-\Pj_\bZ)\bW
  =\begin{pmatrix}z^T(I-\Pj_\bZ)z & z^T(I-\Pj_\bZ)w\\ w^T(I-\Pj_\bZ)z & w^T(I-\Pj_\bZ)w\end{pmatrix}
  =\begin{pmatrix}S_{zz}&S_{zw}\\S_{zw}&S_{ww}\end{pmatrix}=:M,
\]
\[
  \bW^T(I-\Pj_\bZ)\bY=\begin{pmatrix}z^T(I-\Pj_\bZ)\bY\\ w^T(I-\Pj_\bZ)\bY\end{pmatrix}
  =\begin{pmatrix}S_{zy}\\S_{wy}\end{pmatrix},
\]
which gives $\wh\bb=M^{-1}(S_{zy},S_{wy})^T$. \qed \ \ Also $\Var(\wh\bb)=\sigma^2M^{-1}$.

\pt{b} This is $H_0:L^T\bb_{\text{full}}=0$ with a single constraint ($r=1$) picking out
$\beta_2$. The likelihood ratio test is equivalent to the $F$/Wald test (reject for small
$\Lambda=(\SSE_{full}/\SSE_{red})^{n/2}$ $\iff$ reject for large
$(\SSE_{red}-\SSE_{full})/\SSE_{full}$ $\iff$ reject for large $F$), and with one constraint
\[
  F=W=\frac{\wh\beta_2^{\,2}}{\wh\sigma^2\,[M^{-1}]_{22}} .
\]
Now invert the $2\times2$ matrix:
\[
  M^{-1}=\frac{1}{S_{zz}S_{ww}-S_{zw}^2}\begin{pmatrix}S_{ww}&-S_{zw}\\-S_{zw}&S_{zz}\end{pmatrix}
  \quad\Longrightarrow\quad
  [M^{-1}]_{22}=\frac{S_{zz}}{S_{zz}S_{ww}-S_{zw}^2}=\frac{1}{S_{ww}-\dfrac{S_{zw}^2}{S_{zz}}} .
\]
Therefore
\[
  \boxed{\ F=\Big(S_{ww}-\frac{S_{zw}^2}{S_{zz}}\Big)\frac{\wh\beta_2^{\,2}}{\wh\sigma^2}=U
  \ \overset{H_0}{\sim}\ F_{1,\,n-r-2}\ }
\]
and the LRT rejects $H_0$ for large $U$. \qed

\begin{idea}{Reading the coefficient}
$S_{ww}-S_{zw}^2/S_{zz}=w^T(I-\Pj_{[\bZ:z]})w$ is the residual sum of squares of $w$ after
removing \emph{both} the treatment means and $z$. So $U$ is exactly
$\wh\beta_2^{\,2}\times(\text{information about }\beta_2)/\wh\sigma^2$: the more of $w$ that
survives adjustment for $z$, the more powerful the test --- the same collinearity phenomenon
as 2024 Q4(c).
\end{idea}

\pt{c} Now $r=2$ constraints, $L^T\bb_{\text{full}}=\bb=(\beta_1,\beta_2)^T=\zero$. The Wald
statistic is
\[
  W=(\wh\bb-\zero)^T\big(\widehat{\Var}(\wh\bb)\big)^{-1}(\wh\bb-\zero)
   =\wh\bb^T\big(\wh\sigma^2M^{-1}\big)^{-1}\wh\bb
   =\frac{1}{\wh\sigma^2}\,\wh\bb^TM\,\wh\bb=V,
\]
and the LRT rejects for large $V$, with
\[
  \boxed{\ \frac{V}{2}\ \overset{H_0}{\sim}\ F_{2,\,n-r-2}. }
\]
Equivalently $V=(\SSE_{red}-\SSE_{full})/\wh\sigma^2$, where $\SSE_{red}$ is the residual sum
of squares of the pure one-way ANOVA model $Y_{ij}=\mu_i+\varepsilon_{ij}$. \qed

%------------------------------------------------------------------
\qhead{8}{}
\begin{stmt}
$\mathbf{A}$ and $\mathbf{B}$ are symmetric $n\times n$ matrices, $\mathbf{B}$ is idempotent,
and $\mathbf{A}$ and $I_n-\mathbf{A}-\mathbf{B}$ are non-negative definite. Show
$\mathbf{AB}=\mathbf{BA}=\mathbf{O}_n$.
\emph{[Hint: show $y^T\mathbf{A}y=0$ for all $y\in\Cs(\mathbf{B})$.]}
\end{stmt}

\sol \textbf{Step 1: $\mathbf{B}$ acts as the identity on its own column space.}
$\mathbf{B}$ is symmetric and idempotent, hence the orthogonal projector onto
$\Cs(\mathbf{B})$. So for every $y\in\Cs(\mathbf{B})$, writing $y=\mathbf{B}u$,
\[
  \mathbf{B}y=\mathbf{B}^2u=\mathbf{B}u=y
  \qquad\Longrightarrow\qquad y^T\mathbf{B}y=y^Ty .
\]

\textbf{Step 2: squeeze $y^T\mathbf{A}y$ between $0$ and $0$.} Let $y\in\Cs(\mathbf{B})$.
Non-negative definiteness of $I-\mathbf{A}-\mathbf{B}$ gives
\[
  0\ \le\ y^T(I-\mathbf{A}-\mathbf{B})y
   \ =\ y^Ty-y^T\mathbf{A}y-\underbrace{y^T\mathbf{B}y}_{=\,y^Ty}
   \ =\ -\,y^T\mathbf{A}y,
\]
so $y^T\mathbf{A}y\le0$. But $\mathbf{A}\succeq0$ gives $y^T\mathbf{A}y\ge0$. Hence
\[
  y^T\mathbf{A}y=0\qquad\text{for every } y\in\Cs(\mathbf{B}).
\]

\textbf{Step 3: from the quadratic form to the matrix.} Since $\mathbf{A}$ is symmetric and
non-negative definite it has a square root: $\mathbf{A}=L^TL$ (take
$L=\mathbf{A}^{1/2}=P\Lambda^{1/2}P^T$ from the spectral decomposition). Then for
$y\in\Cs(\mathbf{B})$,
\[
  0=y^T\mathbf{A}y=y^TL^TLy=\|Ly\|^2
  \quad\Longrightarrow\quad Ly=\zero
  \quad\Longrightarrow\quad \mathbf{A}y=L^TLy=\zero .
\]

\textbf{Step 4: conclude.} Every column of $\mathbf{B}$ lies in $\Cs(\mathbf{B})$, so
$\mathbf{A}$ annihilates each of them:
\[
  \mathbf{AB}=\mathbf{O}.
\]
Transposing and using symmetry of both matrices,
$\mathbf{O}=(\mathbf{AB})^T=\mathbf{B}^T\mathbf{A}^T=\mathbf{BA}$. \qed

\begin{confusion}{``$y^T\mathbf{A}y=0$'' does not by itself give $\mathbf{A}y=\zero$}
For a general symmetric $\mathbf{A}$ it does not --- e.g.\
$\mathbf{A}=\begin{bsmallmatrix}0&1\\1&0\end{bsmallmatrix}$ and $y=e_1$ give
$y^T\mathbf{A}y=0$ but $\mathbf{A}y=e_2\ne\zero$. Step 3 needs
$\mathbf{A}\succeq0$, which is exactly why that hypothesis is in the question. State this.
\end{confusion}

%==================================================================
\section{Practice: linear-algebra proofs at paper level}

Five questions in the style of 2024 Q1/Q3 and 2025 Q4 --- pure matrix algebra, no data.
Q\,3.1 is Homework~2 Q4 of this year, so it is the likeliest of the five.

\qhead{10}{}
\begin{stmt}
$A$ is an $n\times n$ positive semidefinite matrix with $\rank(A)=r$ and
$\rank(I_n-A)=n-r$, for some $1\le r<n$. Prove that $A$ is an orthogonal projection matrix.
In addition, if $\bY\sim\N(\bmu,I_n)$ and $U_1:=\bY^TA\bY$, $U_2:=\bY^T\bY-\bY^TA\bY$, show
--- \emph{without using the Fisher--Cochran theorem} --- that $U_1$ and $U_2$ are
independent, with $U_1\sim\chi^2_r(\lambda_1/2)$ and $U_2\sim\chi^2_{n-r}(\lambda_2/2)$,
where $\lambda_1=\bmu^TA\bmu$ and $\lambda_2=\bmu^T\bmu-\bmu^TA\bmu$.
\end{stmt}

\pt{Part 1: $A$ is an orthogonal projector} $A$ is positive semidefinite, hence symmetric, so
spectrally decompose
\[
  A=P\Lambda P^T,\qquad P^TP=PP^T=I_n,\qquad \Lambda=\mathrm{diag}(\lambda_1,\dots,\lambda_n),\ \lambda_i\ge0 .
\]
\begin{itemize}
  \item $\rank(A)=r$ $\Longrightarrow$ \textbf{exactly $r$} of the $\lambda_i$ are non-zero.
  \item $I_n-A=P(I-\Lambda)P^T$ has eigenvalues $1-\lambda_i$, so $\rank(I_n-A)=n-r$
        $\Longrightarrow$ exactly $n-r$ of the $1-\lambda_i$ are non-zero $\Longrightarrow$
        \textbf{exactly $r$} of the $\lambda_i$ equal $1$.
\end{itemize}
Those $r$ eigenvalues equal to $1$ are certainly non-zero, and there are only $r$ non-zero
ones altogether. Hence the non-zero eigenvalues are \emph{precisely} those $r$ ones equal to
$1$:
\[
  \lambda_i\in\{0,1\}\ \ \forall i,\qquad\text{with exactly } r \text{ ones.}
\]
Therefore $\Lambda^2=\Lambda$, so $A^2=P\Lambda^2P^T=P\Lambda P^T=A$. Being symmetric and
idempotent, $A$ is an orthogonal projection matrix. \qed

\pt{Part 2: distributions and independence} Split
$P=[\,C:D\,]$, where $C$ is the $n\times r$ block of eigenvectors with eigenvalue $1$ and $D$
the $n\times(n-r)$ block with eigenvalue $0$. Orthogonality of $P$ gives
\[
  C^TC=I_r,\qquad D^TD=I_{n-r},\qquad C^TD=0,\qquad A=CC^T,\qquad I-A=DD^T .
\]
Hence
\[
  U_1=\bY^TCC^T\bY=\|C^T\bY\|^2,
  \qquad
  U_2=\bY^T(I-A)\bY=\|D^T\bY\|^2 .
\]
Put $Z_1=C^T\bY$ and $Z_2=D^T\bY$. As linear images of a normal vector these are jointly
normal, with
\[
  Z_1\sim\N_r(C^T\bmu,\ C^TC)=\N_r(C^T\bmu,I_r),
  \qquad
  Z_2\sim\N_{n-r}(D^T\bmu,I_{n-r}),
\]
\[
  \Cov(Z_1,Z_2)=C^TI_nD=C^TD=0 .
\]
Jointly normal and uncorrelated $\Longrightarrow$ independent; therefore
$U_1=\|Z_1\|^2$ and $U_2=\|Z_2\|^2$ are independent. Finally, a squared norm of a
$\N_k(\theta,I_k)$ vector is $\chi^2_k(\|\theta\|^2/2)$, and
\[
  \|C^T\bmu\|^2=\bmu^TCC^T\bmu=\bmu^TA\bmu=\lambda_1,
  \qquad
  \|D^T\bmu\|^2=\bmu^T(I-A)\bmu=\bmu^T\bmu-\bmu^TA\bmu=\lambda_2 ,
\]
so $U_1\sim\chi^2_r(\lambda_1/2)$ and $U_2\sim\chi^2_{n-r}(\lambda_2/2)$. \qed

\begin{idea}{Why ``without Fisher--Cochran'' is not a handicap}
Once $A$ is known to be a projector you can \emph{split the space}: $C$ and $D$ are two
blocks of a single orthogonal matrix, so independence is just $C^TD=0$. Fisher--Cochran is
the general $r$-piece theorem; this is the two-piece case done by hand, and it is shorter.
\end{idea}

%------------------------------------------------------------------
\qhead{8}{}
\begin{stmt}
Show that $G$ is a generalised inverse of $A$ if and only if $AG$ is idempotent and
$\rank(AG)=\rank(A)$.
\end{stmt}

\pt{$\Rightarrow$} Suppose $AGA=A$. Then
\[
  (AG)^2=AGAG=(AGA)G=AG,
\]
so $AG$ is idempotent. For the rank, $\rank(AG)\le\rank(A)$ always, while
$\rank(A)=\rank(AGA)\le\rank(AG)$. Hence $\rank(AG)=\rank(A)$.

\pt{$\Leftarrow$} Suppose $AG$ is idempotent with $\rank(AG)=\rank(A)$. Since
$\Cs(AG)\subseteq\Cs(A)$ and the two dimensions agree,
\[
  \Cs(AG)=\Cs(A).
\]
An idempotent matrix acts as the identity on its own column space: if $u\in\Cs(AG)$, say
$u=AGw$, then $(AG)u=(AG)^2w=AGw=u$. Now take any $x$. Then $Ax\in\Cs(A)=\Cs(AG)$, so
$Ax=AGu$ for some $u$, and
\[
  AGAx=AG(AGu)=(AG)^2u=AGu=Ax .
\]
This holds for every $x$, so $AGA=A$. \qed

\begin{confusion}{the rank condition cannot be dropped}
$AG$ idempotent alone is not enough: $G=O$ makes $AG=O$ idempotent, but $AGA=O\ne A$. The
rank condition is what forces $\Cs(AG)$ to be all of $\Cs(A)$.
\end{confusion}

%------------------------------------------------------------------
\qhead{6}{}
\begin{stmt}
For an $n\times p$ matrix $\bX$, show that $\Nsp(\bX)=\Nsp(\bX^T\bX)$, and deduce that
$\Cs(\bX^T)=\Cs(\bX^T\bX)$ and $\rank(\bX)=\rank(\bX^T\bX)$.
\end{stmt}

\sol \textbf{$\Nsp(\bX)\subseteq\Nsp(\bX^T\bX)$:} if $\bX v=\zero$ then $\bX^T\bX v=\zero$.

\textbf{$\Nsp(\bX^T\bX)\subseteq\Nsp(\bX)$:} if $\bX^T\bX v=\zero$ then
\[
  0=v^T\bX^T\bX v=\|\bX v\|^2 \quad\Longrightarrow\quad \bX v=\zero .
\]
So the two null spaces coincide. Both $\bX$ and $\bX^T\bX$ have $p$ columns, so rank--nullity
gives
\[
  \rank(\bX)=p-\dim\Nsp(\bX)=p-\dim\Nsp(\bX^T\bX)=\rank(\bX^T\bX).
\]
Finally $\Cs(\bX^T\bX)\subseteq\Cs(\bX^T)$ is immediate, and
\[
  \dim\Cs(\bX^T\bX)=\rank(\bX^T\bX)=\rank(\bX)=\rank(\bX^T)=\dim\Cs(\bX^T),
\]
so a subspace of equal dimension must be the whole space: $\Cs(\bX^T)=\Cs(\bX^T\bX)$. \qed

\begin{recipe}{The two consequences to quote}
\begin{itemize}
  \item The normal equations $\bX^T\bX\bb=\bX^T\bY$ are \textbf{always consistent}, because
        $\bX^T\bY\in\Cs(\bX^T)=\Cs(\bX^T\bX)$.
  \item $a^T\bb$ is estimable $\iff a\in\Cs(\bX^T)=\Cs(\bX^T\bX)$ --- the two forms of the
        criterion are the same statement.
\end{itemize}
\end{recipe}

%------------------------------------------------------------------
\qhead{10}{}
\begin{stmt}
Let $G$ be any generalised inverse of $\bX^T\bX$ and set $P=\bX G\bX^T$. Show that $P$ is the
orthogonal projection matrix onto $\Cs(\bX)$, and that it does not depend on which g-inverse
$G$ is chosen.
\end{stmt}

\pt{Step 1: $\bX G\bX^T\bX=\bX$} Put $M=\bX G\bX^T\bX-\bX$ and expand $M^TM$, using
$\bX^T\bX G\bX^T\bX=\bX^T\bX$ repeatedly:
\begin{align*}
  M^TM &= \bX^T\bX G^T\big(\bX^T\bX G\bX^T\bX\big)-\bX^T\bX G^T\bX^T\bX
         -\big(\bX^T\bX G\bX^T\bX\big)+\bX^T\bX\\
       &= \bX^T\bX G^T\bX^T\bX-\bX^T\bX G^T\bX^T\bX-\bX^T\bX+\bX^T\bX=O .
\end{align*}
The diagonal entries of $M^TM$ are the squared norms of the columns of $M$, so $M^TM=O$
forces $M=O$. Hence $\bX G\bX^T\bX=\bX$, i.e.\ $P\bX=\bX$.

\pt{Step 2: idempotent} $P^2=\bX G\bX^T\cdot\bX G\bX^T=\big(\bX G\bX^T\bX\big)G\bX^T=\bX G\bX^T=P$.

\pt{Step 3: the image is $\Cs(\bX)$} $\Cs(P)\subseteq\Cs(\bX)$ from the leading factor $\bX$;
and $P\bX=\bX$ gives $\Cs(\bX)\subseteq\Cs(P)$. So $\Cs(P)=\Cs(\bX)$.

\pt{Step 4: invariance} Use the rule \emph{$AB^-C$ is invariant to the choice of $B^-$ when
$\mathcal{R}(A)\subseteq\mathcal{R}(B)$ and $\Cs(C)\subseteq\Cs(B)$}, with $A=\bX$,
$B=\bX^T\bX$, $C=\bX^T$. By Q\,3.3,
\[
  \mathcal{R}(\bX)=\Cs(\bX^T)=\Cs(\bX^T\bX)=\mathcal{R}(B),
  \qquad
  \Cs(\bX^T)=\Cs(\bX^T\bX)=\Cs(B),
\]
both with equality. Hence $P=\bX G\bX^T$ is the same matrix for every g-inverse $G$.

\pt{Step 5: symmetry} A \emph{symmetric} g-inverse of $\bX^T\bX$ always exists: spectrally
decompose $\bX^T\bX=P_0\Lambda P_0^T$ and take $G_0=P_0\Delta P_0^T$ with
$\delta_i=1/\lambda_i$ where $\lambda_i\ne0$ and $\delta_i=0$ elsewhere. For that choice
$\bX G_0\bX^T$ is visibly symmetric --- and by Step 4 every choice gives the \emph{same}
matrix. So $P$ is symmetric.

Symmetric, idempotent, with image $\Cs(\bX)$: $P$ is the orthogonal projector onto
$\Cs(\bX)$. \qed

\begin{idea}{The move worth stealing}
Prove \textbf{invariance first}, then verify the awkward property (here, symmetry) for one
conveniently chosen $G$. Trying to show $\bX G\bX^T$ is symmetric for an \emph{arbitrary} $G$
is painful and unnecessary.
\end{idea}

%------------------------------------------------------------------
\qhead{6}{}
\begin{stmt}
$A$ and $B$ are symmetric idempotent $n\times n$ matrices with $\Cs(B)\subseteq\Cs(A)$. Show
that $AB=BA=B$, that $A-B$ is an orthogonal projection matrix, and that
$\rank(A-B)=\rank(A)-\rank(B)$.
\end{stmt}

\sol \textbf{$AB=B$.} $A$ is symmetric and idempotent, hence the orthogonal projector onto
$\Cs(A)$, so $Au=u$ for every $u\in\Cs(A)$. For any $x$ we have $Bx\in\Cs(B)\subseteq\Cs(A)$,
so $A(Bx)=Bx$. As this holds for all $x$, $AB=B$.

\textbf{$BA=B$.} Transpose, using symmetry of both: $BA=B^TA^T=(AB)^T=B^T=B$.

\textbf{$A-B$ is an orthogonal projector.} It is symmetric as a difference of symmetric
matrices, and
\[
  (A-B)^2=A^2-AB-BA+B^2=A-B-B+B=A-B .
\]

\textbf{Rank.} $A-B$ is idempotent, so trace counts rank:
\[
  \rank(A-B)=\tr(A-B)=\tr(A)-\tr(B)=\rank(A)-\rank(B). \qquad\blacksquare
\]

\begin{idea}{Where you have already used this}
This is exactly why $P_\bX-P_\bZ$ is a projector of rank $\rank(\bX)-\rank(\bZ)$ in the
nested-model decomposition $P_\bX=P_\bZ+P_{\bW\mid\bZ}$ --- and hence why
$\mathrm{SSE}_{red}-\mathrm{SSE}_{full}$ is a chi-square with the right degrees of freedom.
Every ANOVA table rests on it.
\end{idea}

\begin{recipe}{The five moves these questions keep testing}
\begin{enumerate}
  \item Eigenvalues of an idempotent matrix lie in $\{0,1\}$, hence $\tr=\rank$.
        \hfill(Q\,3.1, Q\,3.5)
  \item Spectral decomposition to split $\R^n$ into orthogonal blocks. \hfill(Q\,3.1)
  \item $M^TM=O\Rightarrow M=O$, and $v^TA v=\|\bX v\|^2=0\Rightarrow \bX v=\zero$.
        \hfill(Q\,3.3, Q\,3.4)
  \item An idempotent matrix acts as the identity on its own column space.
        \hfill(Q\,3.2, Q\,3.5)
  \item Prove invariance first, then check the property for one convenient choice.
        \hfill(Q\,3.4)
\end{enumerate}
\textbf{If a question asks you to prove anything about projections, the first line is
``symmetric and idempotent'' and the second is $\tr=\rank$.}
\end{recipe}

%==================================================================
\section{Coverage audit: do the crash courses cover all of this?}

\textbf{They do now.} When these solutions were first written, the machinery for every
question was in the crash courses but four specific results that were actually asked were
not. Those four have since been \textbf{patched into Crash Courses 2 and 3}, and the proof of
the third fundamental theorem has been filled in. The table below records where each one now
lives.

\begin{center}
\begin{tabular}{@{}llp{6.4cm}l@{}}
\toprule
Paper & Q & Where the method is & Status \\
\midrule
2024 & 1 & CC2 \S2.4 \emph{Sums and differences of projections} & \textbf{Added} \\
2024 & 2 & CC2 \S4--5 (estimability, G--M); CC3 \S6 (test) & Already covered \\
2024 & 3 & CC2 \S2.8 \emph{Building a g-inverse from a full-rank corner} & \textbf{Added} \\
2024 & 4 & CC3 \S2.4(a),(b) \emph{Two consequences} & \textbf{Added} \\
2025 & 1 & CC3 \S7 (3rd fundamental theorem, proof now given) & \textbf{Strengthened} \\
2025 & 2 & CC2 \S4.2 (the $2\times2$ $\alpha_i+\gamma_j$ example) & Already covered \\
2025 & 3 & CC3 \S2.4(c) \emph{Analysis of covariance} & \textbf{Added} \\
2025 & 4 & CC2 \S2.5 \emph{A non-negative-definite squeeze} & \textbf{Added} \\
\bottomrule
\end{tabular}
\end{center}

\begin{recipe}{The four results that were missing, and where they now are}
\begin{enumerate}
  \item \textbf{$\Pj+\mathbf{Q}$ is a projector $\iff\Pj\mathbf{Q}=\mathbf{O}$}, with rank
        additivity (2024 Q1). Now \textbf{CC2 \S2.4}, together with the mirror result that
        $\Pj-\mathbf{Q}$ is a projector when $\Cs(\mathbf{Q})\subseteq\Cs(\Pj)$ --- which is
        what makes $P_\bX-P_\bZ$ a projector in the nested-model decomposition.
  \item \textbf{The block g-inverse} $G=\begin{bsmallmatrix}A_{11}^{-1}&0\\0&0\end{bsmallmatrix}$
        (2024 Q3). Now \textbf{CC2 \S2.8}, with the point that the rank hypothesis is used in
        exactly one place: forcing the Schur complement to vanish.
  \item \textbf{$\mathbf{A}\succeq0$, $I-\mathbf{A}-\mathbf{B}\succeq0$,
        $\mathbf{B}^2=\mathbf{B}\Rightarrow\mathbf{AB}=\mathbf{BA}=\mathbf{O}$} (2025 Q4).
        Now \textbf{CC2 \S2.5}, with the counterexample showing why $\mathbf{A}\succeq0$
        cannot be dropped.
  \item \textbf{$\Var(\wh\beta_j)=\sigma^2/d_j$, the $T_j$ decomposition and the collinearity
        argument} (2024 Q4), and the \textbf{ANCOVA} instance of Frisch--Waugh with the
        $S_{zz},S_{zw},\dots$ notation and the $U$, $V$ statistics (2025 Q3). Now
        \textbf{CC3 \S2.4}, parts (a)--(c).
\end{enumerate}
\end{recipe}

\begin{idea}{What the two papers never asked}
No numerical ANOVA table, no calculator, no confidence intervals, no prediction, no $R^2$.
Both papers were pure theory. \textbf{This year's guide breaks that pattern} --- it promises
a numerical problem and three hours instead of two --- so Crash Course 4 (one-way ANOVA) and
\S10 of the formula sheet cover ground the past papers do not test. Prepare both.
\end{idea}

\vfill
\begin{center}\small\itshape
Companions: the four crash courses, and the two-page formula sheet.
\end{center}

\end{document}
